\introduction

\paragraph*{Актуальность темы исследования.}

Электронная коммерция продолжает расти. Годовой объём российского сегмента в 2025 году составил 11,5~трлн~руб. \cite{ref_1}, мировой рынок тоже показывает рост \cite{ref_2}. Каталоги крупных интернет-магазинов содержат десятки миллионов товарных позиций: только индексируемая часть каталога Ozon в открытых аналитических базах превышает 38~млн позиций \cite{ref_3}. Ежедневно туда добавляются тысячи новых записей от партнёров-поставщиков.

Карточка товара — центральный элемент цифровой витрины. От её полноты зависят фасетный поиск, релевантность ранжирования и рекомендаций, корректность аналитики и в итоге конверсия в покупки. Партнёрский ввод при этом неоднородный: для одного товара поставщик передаёт только название и бренд, для другого — расширенный набор полей с фотографиями, для третьего — состав на иностранном языке. Образующиеся пробелы исторически закрывают контент-менеджеры вручную: каждую карточку приходится открывать, проверять и дописывать десятки атрибутов. На каталоге в миллионы позиций и тысячах новых записей в сутки это уже не работает ни по затратам, ни по доступности знаний иностранных языков \cite{ref_4}.

Системы управления товарной информацией (Product Information Management) — Akeneo, Pimcore, Salsify и аналогичные \cite{ref_5, ref_6, ref_7} — сложились из задачи хранить уже сформированную товарную информацию. Задачу извлечения этой информации из неструктурированного партнёрского ввода они системно не решают. Появляющиеся в них модули автоматического обогащения — это, как правило, обёртка над одной публичной большой языковой моделью с прямым вызовом на каждый товар. Подробное сопоставление приведено в~\ref{subsubsec:analogs}. Такой вызов экономически приемлем только для каталогов небольшого и среднего размера и не даёт гарантий фактической корректности. На числовых полях большие языковые модели регулярно ошибаются (\ref{sec:failure-taxonomy}). Рост ставок коммерческих API, ужесточение требований к локализации данных и отраслевой тренд на открытые модели в локальном развёртывании \cite{ref_8} дополнительно усиливают запрос на архитектурные решения, которые не сводятся к прямому вызову одной модели.

Отсюда вытекает противоречие, которое и определяет актуальность работы. С одной стороны, автоматизация наполнения карточек — инженерная необходимость. С другой стороны, ни одна отдельная технология не покрывает все типы целевых атрибутов с приемлемой точностью и стоимостью одновременно. Правила надёжны на числовых атрибутах, но плохо покрывают поля со сложным смыслом. Классические классификаторы требуют большой размеченной выборки. Большие языковые модели универсальны, но дорогие и часто ошибаются на числовых полях.

В качестве разрешения этого противоречия в работе рассматривается гибридная каскадная архитектура. Это последовательность слоёв, в которой каждый следующий подключается только к тому остатку, который не закрыли предыдущие с нужной уверенностью. Такая композиция совмещает семантические возможности генерирующих моделей с предсказательной силой классического машинного обучения и автоматическую отбраковку недостоверных значений. С научной стороны актуальность работы — это проверка эффективности гибридного каскада на типовой задаче обогащения. С практической стороны — создание готового к производству комплекса с контролируемой стоимостью и возможностью локального развёртывания.

\paragraph*{Цель работы.}

Цель работы — разработать гибридную каскадную систему обогащения товарных данных интернет-магазина. Система должна точно заполнять категориальные атрибуты, обходиться дешевле прямого вызова большой языковой модели и выигрывать у прямого вызова одновременно и по точности, и по стоимости.

\paragraph*{Задачи работы.}

Для достижения поставленной цели сформулированы и решены пять задач:

\begin{enumerate}[label=\arabic*)]
    \item разработать архитектуру гибридного каскадного конвейера обработки товарных данных с пошаговой фильтрацией по уверенности и явным разделением слоёв. Слои каскада: предкаскадный маршрутизатор категорий, экстрактор регулярных выражений, классификатор машинного обучения на векторных представлениях, байесовский валидатор плотностей, запасной слой большой языковой модели. Пошаговая фильтрация формализует требование утверждённой темы об ансамблевой оценке уверенности моделей как последовательное конъюнктивное условие на калиброванные пороги слоёв;
    \item реализовать каждый слой каскада в виде отдельного программного компонента: экстрактор регулярных выражений; классификатор машинного обучения на многоязычных векторных представлениях партнёр-доступных полей; байесовский валидатор, который использует обученную сеть атрибутов как фильтр статистически маловероятных пар «бренд × значение»; запасной слой на основе большой языковой модели с поддержкой нескольких поставщиков, включая открытые модели для локального развёртывания;
    \item построить эталонную разметку и проверочную выборку без пересечения товарных кодов для оценки качества системы на трёх категориях: паста, шоколад, сыры. Эталон сочетает слабую разметку из тегов открытого каталога, поатрибутное согласие трёх независимых больших языковых моделей и слепую разметку референсной моделью;
    \item провести экспериментальную оценку разработанной системы: измерить точность сквозного прохождения каскада, поатрибутное покрытие и стоимость обработки на эталонной тестовой выборке; сопоставить полученные показатели с прямым использованием большой языковой модели; оценить вклад каждого слоя каскада, устойчивость к смещению по языкам входных данных и поведение байесовского валидатора как селективной валидации плотностей;
    \item разработать демонстрационный программный комплекс. Комплекс должен обеспечивать работу оператора каталога с каскадом через веб-интерфейс и стабильный программный интерфейс. Поддерживаются три производственных сценария развёртывания: с премиум-коммерческой большой языковой моделью, с бюджетной коммерческой и с открытой моделью в локальном развёртывании.
\end{enumerate}

Каждой задаче соответствует измеримый результат, изложенный в подразделе «Новизна и основные результаты».

\paragraph*{Объект и предмет исследования.}

\textbf{Объектом исследования} является процесс автоматизированного обогащения товарных каталогов интернет-магазинов на основе открытых источников данных. На вход подаётся минимальный набор партнёр-доступных полей, на выходе требуется получить структурированный набор целевых атрибутов.

\textbf{Предметом исследования} является гибридная каскадная архитектура обработки товарных данных с ансамблевой оценкой уверенности классификаторов и выборочным обращением к большой языковой модели в роли запасного слоя. Рассматриваются её точность, стоимость, устойчивость к смещению входного распределения, а также вычислительные и эксплуатационные свойства, существенные для производственного применения.

\paragraph*{Теоретические и методические основы работы.}

Работа опирается на следующие методы и инструменты:

\begin{enumerate}[label=\arabic*),leftmargin=*]
    \item извлечение признаков из неструктурированного текста --- регулярные выражения и многоязычный кодировщик \texttt{paraphrase-\allowbreak multilingual-\allowbreak mpnet-\allowbreak base-v2};
    \item классификатор уровня машинного обучения --- градиентный бустинг XGBoost поверх векторных представлений;
    \item валидатор атрибутов реализован вероятностной графической моделью на библиотеке pgmpy; структура сети строится алгоритмом восхождения к вершине (англ. Hill Climbing) по байесовскому информационному критерию (BIC);
    \item обучающая выборка формируется слабой разметкой по тегам открытого каталога Open Food Facts;
    \item эталонная разметка строится из трёх ярусов: детерминированные правила для нутриент-производных атрибутов, регуляторные теги и поатрибутное согласие трёх независимых больших языковых моделей со слепой проверкой референсной моделью; процедура подробно описана в~\ref{subsubsec:gold_labeling};
    \item качество измеряется на отложенной выборке без пересечения товарных кодов; доверительные интервалы --- кластерным бутстрэпом с группировкой по товарным кодам;
    \item гипотезы о превосходстве каскада над одиночной языковой моделью проверяются по заранее объявленному протоколу с поправкой Бонферрони на множественные тесты.
\end{enumerate}

Архитектурный подход к выбору модели по соотношению цена/качество опирается на работы L. Chen и соавторов FrugalGPT \cite{ref_9}, P. Aggarwal и соавторов AutoMix \cite{ref_10}, D. Ding и соавторов Hybrid LLM \cite{ref_11}, а также на бенчмарк Q. J. Hu и соавторов RouterBench \cite{ref_12}.

\paragraph*{Новизна и основные результаты.}

Каждой из сформулированных выше задач в работе соответствует не менее одного измеримого результата.

\textbf{По задаче 1} — проектирование архитектуры каскада, разделы~\ref{subsec:formal_task}, \ref{subsubsec:layers_impl}. Спроектирован гибридный каскад из пяти последовательных слоёв; решение слоя принимается при одновременном выполнении калиброванных порогов уверенности предыдущих слоёв. Формальная постановка задачи допускает явный отказ системы от ответа и явно указывает слой-источник каждого предсказания.

\textbf{По задаче 2} — реализация слоёв, раздел~\ref{subsubsec:demo_api}. Каскад реализован как набор независимо обновляемых программных компонентов. Разработана демонстрационная система: шлюз API на Go — ML-сервис на Python/FastAPI — статический веб-интерфейс. Система запускается одной командой и проходит проверку готовности всех слоёв через диагностический интерфейс.

\textbf{По задаче 3} — построение эталонной разметки, раздел~\ref{subsec:data}. Построена многоисточниковая эталонная разметка для трёх категорий товаров: теги открытого каталога, правила на нутриентах, поатрибутное согласие трёх независимых больших языковых моделей и слепая разметка референсной моделью. Разбиение на обучающую, валидационную и тестовую части выполнено без пересечения товарных кодов.

\textbf{По задаче 4} — экспериментальная оценка, разделы~\ref{subsec:headline_result}, \ref{subsec:layer_accuracy}, \ref{subsubsec:bayes_behavior}, \ref{subsec:per_language}. Получена сквозная точность 93,0~\% (точность только-каскадной части 95,5~\%, macro-F1 0,890) на тестовой выборке без пересечения товарных кодов (новые товарные позиции) из 16\,360 ячеек по 20 парам «категория-атрибут» из 21 атрибута производственной схемы при 580-кратном сокращении стоимости в режиме «каскад + Gemini 2.5 Flash» относительно прямого вызова Claude Sonnet 4.5 (сюда входит архитектурный вклад каскада в 24 раза и удешевление модели). На независимом эталоне ручной разметки Claude Opus 4 (n=566 ячеек с ответом каскада из 688 в области применимости) сквозная точность — 86,7~\%. Выигрыш каскада по точности и стоимости одновременно подтверждён на пяти моделях запасного слоя. Вклад каждого слоя измерен количественно, включая точечный вклад байесовского валидатора. Каскад покрывает 97,3~\% ячеек.

\textbf{По задаче 5} — разработка демонстрационного программного комплекса, разделы~\ref{subsec:user_guide}, \ref{subsec:scenarios}. Реализована поддержка трёх производственных сценариев развёртывания: премиум-коммерческий, бюджетный коммерческий, локальный с открытой моделью. Поддерживаются три роли пользователей: контент-менеджер, системный аналитик, инженер машинного обучения.

На основании результатов сформулированы три пункта научной новизны. Их общая особенность относительно известных стоимостно-ориентированных каскадов для языковых задач — FrugalGPT \cite{ref_9}, AutoMix \cite{ref_10}, Hybrid LLM \cite{ref_11}, RouterBench \cite{ref_12} — разнородность слоёв. Перечисленные работы комбинируют однотипные генеративные модели разной мощности; в настоящей работе исследован разнородный каскад «правила (регулярные выражения) — дискриминативная модель машинного обучения — генеративная большая языковая модель» с явным правилом маршрутизации на уровне пары «категория-атрибут». Подобной комбинации трёх парадигм для задачи обогащения товарных каталогов в литературе не зафиксировано.

\textbf{Первый пункт научной новизны.} Впервые количественно установлено доминирование по Парето разработанной гибридной каскадной системы над прямым использованием большой языковой модели на задаче обогащения товарных каталогов. На тестовой выборке без пересечения товарных кодов (новые товарные позиции) размером 16\,360 ячеек одновременно достигаются: превышение точности на 9,2~п.~п. (93,0~\% против 83,8~\%) и сокращение стоимости обработки в 580 раз (сюда входит архитектурный вклад каскада в 24 раза и удешевление модели запасного слоя). Компромисса между точностью и стоимостью при этом не возникает.

\textbf{Второй пункт научной новизны.} Эмпирически установлено: в сценарии локального развёртывания с открытой большой языковой моделью gpt-oss-120b в роли запасного слоя каскад компенсирует слабость малой открытой модели на сложных атрибутах за счёт специализированного слоя машинного обучения. Прирост сквозной точности — 21,3~п.~п. над прямым использованием той же открытой модели (90,6~\% против 69,3~\%) при сокращении стоимости в 810 раз относительно опорного прямого вызова Claude Sonnet 4.5. То есть архитектура инвариантна к выбору поставщика большой языковой модели.

\textbf{Третий пункт научной новизны} объединяет два самостоятельных результата. Первый — заранее объявленный отрицательный исход проверки гипотезы H1 о том, что обучаемый стоимостно-чувствительный маршрутизатор превосходит статическое правило при равном бюджете запасного слоя: впервые на задаче обогащения товарных каталогов по заранее объявленной процедуре с поправкой Бонферрони на множественные тесты эмпирически установлено, что статическое правило маршрутизации на уровне пары (категория, атрибут) не уступает обучаемому XGBoost-маршрутизатору с учётом стоимости по сквозной точности при равных бюджетах вызовов запасного слоя. Результат уточняет применимость обучаемой маршрутизации к этому классу задач и согласуется с выводами бенчмарка RouterBench \cite{ref_12}. Второй — перенос замкнутого цикла «аудит эталонной разметки — правка экстрактора признаков — переобучение моделей — повторное измерение» на категории внутри продуктового домена. Цикл воспроизведён на трёх категориях из трёх с приростом сквозной точности на подвыборке проверенных проблемных ячеек +18,2, +16,3 и +14,4~п.~п. для пасты, шоколада и сыров соответственно.

\paragraph*{Апробация и внедрение результатов.}

Разработка велась по практическим требованиям, типичным для крупных интернет-магазинов (российских и международных). Архитектура и эксплуатационные характеристики системы ориентированы на встраивание в производственный процесс обогащения товарного каталога без перестройки соседних звеньев (см.~\ref{subsubsec:integration} — описание точки встраивания в маршрут «партнёрский поток данных $\to$ каталог $\to$ витрина»). Поддерживаются три производственных сценария развёртывания: с премиум-коммерческой большой языковой моделью (Claude Sonnet 4.5 / GPT-4o), с бюджетной коммерческой (Gemini 2.5 Flash) и с открытой моделью в полностью локальном развёртывании (gpt-oss-120b). Последний сценарий критичен для требований локализации обработки партнёрских данных.

Апробация выполнена в виде демонстрационного программного комплекса. Комплекс запускается одной командой и проверяется на работоспособность через диагностический программный интерфейс. Демонстрация подтверждает, что все слои каскада работают на трёх категориях: паста, шоколад, сыры. Воспроизводимость численных показателей обеспечена сценарием \texttt{reproduce.sh}, который пересоздаёт основные результаты из артефактов в каталоге \texttt{datasets/processed/} репозитория. Разработанный комплекс готов к интеграции в системы управления товарной информацией розничной торговли (см.~\ref{subsubsec:integration}).

\paragraph*{Структура и объём работы.}

Выпускная квалификационная работа состоит из введения, четырёх глав, заключения и списка использованных источников. Общий объём работы --- 115 страниц, основная часть (от введения до заключения включительно) --- 90 страниц. Список использованных источников включает 58 наименований: 18 российских источников, индексируемых в РИНЦ и доступных в открытых российских репозиториях, и 40 зарубежных публикаций, отраслевых отчётов и стандартов.

Первая глава содержит обоснование актуальности темы, литературный обзор по трём направлениям и техническое задание (ТЗ) на разрабатываемую систему. Вторая глава посвящена теоретическому обоснованию решения. В ней приводятся формальная постановка задачи, обоснование архитектуры и стека технологий, а также сценарии прохождения запроса через каскад. Третья глава содержит описание программной реализации, использованных данных и доказательство работоспособности каскада на эталонной тестовой выборке без пересечения товарных кодов. Четвёртая глава представляет порядок применения системы, сценарии её развёртывания и характеристику достигнутых результатов в формате «повышено / сокращено / обеспечено». В заключении собраны выводы по главам, научная новизна и перспективы развития темы.
